
\documentclass[UTF8,zihao=-4]{ctexart}
\usepackage[a4paper,margin=2.15cm]{geometry}
\usepackage{amsmath,amssymb,bm,mathtools}
\usepackage{esint}
\usepackage{booktabs,longtable,tabularx,array,multicol}
\usepackage[most]{tcolorbox}
\usepackage{xcolor}
\usepackage{enumitem}
\usepackage{hyperref}
\usepackage{fancyhdr}
\usepackage{titlesec}
\usepackage{lastpage}
\usepackage{fvextra}
\usepackage{tikz}
\usepackage{eso-pic}
\usepackage{fontspec}
\setCJKmainfont{Noto Serif CJK SC}
\setCJKsansfont{Noto Sans CJK SC}
\setCJKmonofont{Noto Sans Mono CJK SC}
\setmainfont{DejaVu Serif}
\setsansfont{DejaVu Sans}
\setmonofont{DejaVu Sans Mono}
\hypersetup{colorlinks=true,linkcolor=blue!55!black,urlcolor=blue!55!black,citecolor=blue!55!black,pdftitle={面向对象程序设计考前复习知识点讲义},pdfauthor={Jerry}}
\definecolor{mainblue}{HTML}{1F4E79}
\definecolor{lightblue}{HTML}{EAF3FF}
\definecolor{warnred}{HTML}{B3261E}
\definecolor{lightred}{HTML}{FFF1F1}
\definecolor{darkgreen}{HTML}{1B6B3A}
\definecolor{lightgreen}{HTML}{EEF9F0}
\definecolor{orange}{HTML}{B26A00}
\definecolor{lightorange}{HTML}{FFF7E8}
\definecolor{softgray}{HTML}{777777}
\ctexset{section={format=\Large\bfseries\sffamily\color{mainblue}},subsection={format=\large\bfseries\sffamily\color{mainblue!85!black}},subsubsection={format=\normalsize\bfseries\sffamily\color{mainblue!75!black}}}
\titleformat{\paragraph}[runin]{\bfseries\sffamily\color{mainblue!75!black}}{}{0pt}{}[\quad]
\pagestyle{fancy}
\fancyhf{}
\lhead{面向对象程序设计考前复习知识点讲义}
\rhead{\leftmark}
\cfoot{\thepage/\pageref{LastPage}}
\setlength{\headheight}{15pt}
\setlist[itemize]{leftmargin=2em,itemsep=0.25em,topsep=0.25em}
\setlist[enumerate]{leftmargin=2em,itemsep=0.25em,topsep=0.25em}
\renewcommand{\arraystretch}{1.32}
\allowdisplaybreaks
\newcommand{\dd}{\,\mathrm d}
\newcommand{\ee}{\mathrm e}
\newcommand{\ii}{\mathrm i}
\newcommand{\ve}[1]{\bm{#1}}
\newcommand{\grad}{\nabla}
\newcommand{\epsz}{\varepsilon_0}
\newcommand{\muz}{\mu_0}
\newtcolorbox{keybox}[1][]{breakable,colback=lightblue,colframe=mainblue,arc=2mm,title=\textbf{#1},fonttitle=\bfseries\sffamily,coltitle=white}
\newtcolorbox{warnbox}[1][]{breakable,colback=lightred,colframe=warnred,arc=2mm,title=\textbf{#1},fonttitle=\bfseries\sffamily,coltitle=white}
\newtcolorbox{examplebox}[1][]{breakable,colback=lightgreen,colframe=darkgreen,arc=2mm,title=\textbf{#1},fonttitle=\bfseries\sffamily,coltitle=white}
\newtcolorbox{methodbox}[1][]{breakable,colback=lightorange,colframe=orange,arc=2mm,title=\textbf{#1},fonttitle=\bfseries\sffamily,coltitle=white}
\DefineVerbatimEnvironment{Code}{Verbatim}{fontsize=\small,breaklines,breakanywhere,frame=single,framesep=2mm}
\definecolor{watermarkgray}{HTML}{777777}
\AddToShipoutPictureBG{%
  \begin{tikzpicture}[remember picture,overlay]
    \node[rotate=35,scale=1.35,text=watermarkgray,opacity=0.14] at (current page.center) {Jerry 制作｜仅供个人复习使用｜禁止盗印与商用传播};
  \end{tikzpicture}%
}

\begin{document}
\begin{titlepage}
\centering
\vspace*{2.3cm}
{\Huge\bfseries\sffamily 面向对象程序设计考前复习知识点讲义\par}
\vspace{0.75cm}
{\Large 按两套 Java 模拟卷补配的考试导向版\par}
\vspace{0.45cm}
{\large Java 基础 - 类与对象 - 继承多态 - 集合异常 - 程序设计\par}
\vfill
\begin{tcolorbox}[colback=lightblue,colframe=mainblue,width=0.86\textwidth,arc=2mm]
\textbf{使用定位：}这份讲义按已出模拟卷的题型、分值与考点重新整理，重点放在高频结论、题型识别、固定步骤和易错点，不按教材逐段搬运。\par
\textbf{复习方法：}先看题型结构和优先级，再按模块刷小题和大题，最后用公式速查表与最后 30 分钟清单做查漏。
\end{tcolorbox}
\vfill
{\large 资料整理：Jerry\par}
{\large 版本：V1.1\par}
{\large \today\par}
\end{titlepage}
\tableofcontents
\newpage

\section{考试结构与复习策略}

\begin{longtable}{p{0.20\textwidth}p{0.14\textwidth}p{0.15\textwidth}p{0.14\textwidth}p{0.27\textwidth}}
\toprule
\textbf{题型} & \textbf{题量} & \textbf{单题分值} & \textbf{总分} & \textbf{主要考法} \\
\midrule
选择题 & 15 & 2 分 & 30 & Java 基础语法、面向对象概念、常用类 \\
填空题 & 10 空 & 2 分 & 20 & 关键字、方法名、接口、异常、集合 \\
简答题 & 3 & 5 分 & 15 & OOP 三特性、重写、泛型、文件/GUI 事件 \\
程序阅读题 & 3 & 5 分 & 15 & 构造方法、String、继承多态、集合输出 \\
程序设计题 & 1 & 20 分 & 20 & 类封装、ArrayList 管理、增删查改、main 测试 \\
\bottomrule
\end{longtable}

\begin{keybox}[拿分顺序]
选择填空先拿基础关键字与概念；程序阅读题按“对象创建 - 方法调用 - 输出顺序”逐行走；程序设计题先写类、属性、构造、getter/setter，再写 Manager 类和 main 方法，结构完整比炫技更重要。
\end{keybox}

\section{考点优先级}

\begin{longtable}{p{0.16\textwidth}p{0.24\textwidth}p{0.26\textwidth}p{0.24\textwidth}}
\toprule
\textbf{优先级} & \textbf{模块} & \textbf{常见题型} & \textbf{复习要求} \\
\midrule
必考核心 & 类与对象、封装 & 填空、简答、程序设计 & 会写私有属性、构造、getter/setter \\
必考核心 & 继承、多态、接口 & 选择、简答、阅读 & 会区分 extends/implements、重写、向上/向下转型 \\
必考核心 & String、集合、泛型 & 选择、阅读、程序设计 & 会判断 equals、charAt、ArrayList/Set 输出 \\
高频重点 & Java 基础语法 & 选择、填空 & JDK/JRE/JVM、类型转换、数组、标识符 \\
高频重点 & 异常与文件流 & 填空、简答 & try-catch-finally、throws、File、字节流/字符流 \\
可能考点 & 线程与 GUI 事件 & 选择、填空、简答 & start/run、事件源/对象/监听器 \\
了解即可 & 反射、注解 & 选择概念题 & 知道用途，不做复杂代码 \\
\bottomrule
\end{longtable}

\section{模块一：Java 基础语法}

\subsection{JDK、JRE、JVM 与编译运行}
\begin{keybox}[平台无关性]
Java 源文件 \texttt{.java} 经 \texttt{javac} 编译成 \texttt{.class} 字节码，字节码运行在不同平台的 JVM 上。JDK 包含开发工具，JRE 主要用于运行已经编译好的程序。
\end{keybox}

\begin{longtable}{p{0.28\textwidth}p{0.62\textwidth}}
\toprule
\textbf{知识点} & \textbf{考试写法} \\
\midrule
入口方法 & \texttt{public static void main(String[] args)}，方法名是 \texttt{main} \\
源文件命名 & 含 \texttt{public} 类时，源文件名应与 public 类名相同 \\
标识符 & 不能数字开头，不能是关键字，不能含减号，\texttt{\$\_count} 合法 \\
数组长度 & 属性 \texttt{length}，不是 \texttt{length()} \\
类型转换 & 小转大自动拓宽，大转小需强制转换且可能丢失精度 \\
\bottomrule
\end{longtable}

\begin{warnbox}[基础语法易错]
\begin{enumerate}
  \item \texttt{boolean} 不能与 \texttt{int} 互转。
  \item 局部变量使用前必须显式初始化；成员变量有默认值。
  \item \texttt{byte b=128;} 越界；\texttt{double d=10;} 可自动转换。
\end{enumerate}
\end{warnbox}

\section{模块二：类与对象、封装、构造方法}

\subsection{类、对象与构造方法}
\begin{keybox}[核心概念]
\begin{itemize}
  \item 类是对象的模板；对象是类的实例。
  \item \texttt{new} 会在内存中创建对象并调用构造方法。
  \item 构造方法名必须与类名相同，没有返回值类型，可以重载，不能被子类继承。
  \item \texttt{this} 表示当前对象；\texttt{this(...)} 调用本类其他构造方法，且必须是构造方法第一句。
  \item \texttt{super(...)} 调用父类构造方法，也必须是构造方法第一句。
\end{itemize}
\end{keybox}

\begin{methodbox}[封装类书写模板]
\begin{Code}
class Book {
  private String id;
  private String name;

  public Book(String id, String name) {
    this.id = id;
    this.name = name;
  }

  public String getId() { return id; }
  public void setId(String id) { this.id = id; }
}
\end{Code}
\end{methodbox}

\subsection{static、final 与 abstract}
\begin{longtable}{p{0.22\textwidth}p{0.68\textwidth}}
\toprule
\textbf{关键字} & \textbf{常考含义} \\
\midrule
\texttt{static} & 属于类，可通过类名访问；静态方法不能直接访问实例成员 \\
\texttt{final} & 修饰类不能被继承；修饰方法不能被重写；修饰变量赋值后不能再改 \\
\texttt{abstract} & 抽象类不能直接实例化；含抽象方法的类必须声明为抽象类 \\
\bottomrule
\end{longtable}

\section{模块三：继承、多态、接口与类型转换}

\subsection{继承与重写}
\begin{keybox}[继承]
子类继承父类使用 \texttt{extends}。Java 类只支持单继承，但一个类可以实现多个接口。子类可复用父类成员，也可通过重写方法改变行为。
\end{keybox}

\begin{methodbox}[方法重写规则]
\begin{enumerate}
  \item 方法名、参数列表相同。
  \item 返回值类型相同或为协变返回类型。
  \item 访问权限不能比父类更严格。
  \item 不能抛出比父类方法更宽的受检异常。
  \item \texttt{static} 方法不存在运行时多态意义上的重写。
\end{enumerate}
\end{methodbox}

\subsection{多态与类型转换}
\begin{keybox}[多态]
父类引用指向子类对象：
\begin{Code}
Animal a = new Dog();
a.speak();  // 调用子类重写后的 speak
\end{Code}
调用被重写方法时，按实际对象类型动态绑定。父类引用不能直接调用子类新增方法，若要调用，需向下转型并先用 \texttt{instanceof} 判断。
\end{keybox}

\begin{methodbox}[程序阅读：继承多态题]
\begin{enumerate}
  \item 看等号右侧实际对象类型。
  \item 调用重写方法时执行子类版本。
  \item 调用子类特有方法前，必须向下转型。
  \item \texttt{instanceof} 为真再转型，可避免 \texttt{ClassCastException}。
\end{enumerate}
\end{methodbox}

\subsection{接口}
\begin{keybox}[接口]
接口用 \texttt{interface} 定义，类用 \texttt{implements} 实现。接口体现规范。常见考点是“一个类可以实现多个接口”，以及接口引用也可指向实现类对象。
\end{keybox}

\section{模块四：String、包装类、集合与泛型}

\subsection{String 与包装类}
\begin{keybox}[String 常考]
\begin{itemize}
  \item 字符串内容比较用 \texttt{equals()}，不要用 \texttt{==} 判断内容。
  \item \texttt{charAt(0)} 取第一个字符，\texttt{charAt(1)} 取第二个字符。
  \item \texttt{String} 内容不可变；\texttt{StringBuffer} 内容可变。
  \item \texttt{Integer y=10;} 是自动装箱；\texttt{int x=y;} 是自动拆箱。
\end{itemize}
\end{keybox}

\begin{methodbox}[String 输出题模板]
\begin{enumerate}
  \item 字面量拼接如 \texttt{"he"+"llo"} 通常在常量池中得到 \texttt{"hello"}。
  \item \texttt{new String("hello")} 创建新对象。
  \item \texttt{==} 比较引用地址；\texttt{equals()} 比较内容。
\end{enumerate}
\end{methodbox}

\subsection{集合与泛型}
\begin{longtable}{p{0.24\textwidth}p{0.66\textwidth}}
\toprule
\textbf{集合} & \textbf{特点} \\
\midrule
\texttt{ArrayList} & 顺序表，按插入顺序保存，允许重复，适合随机访问 \\
\texttt{LinkedList} & 链表结构，插入删除灵活 \\
\texttt{HashSet} & 不保证顺序，不允许重复 \\
\texttt{TreeSet} & 不允许重复，通常按自然顺序或比较器排序 \\
\texttt{Collection} 遍历 & 迭代器和增强 for 循环 \\
\bottomrule
\end{longtable}

\begin{keybox}[泛型]
泛型用于在编译期约束集合元素类型，减少强制类型转换和运行时类型错误。例如 \texttt{ArrayList<Student>} 表示列表中主要保存 \texttt{Student} 对象。
\end{keybox}

\section{模块五：异常、文件、线程、GUI、注解}

\subsection{异常处理}
\begin{keybox}[异常关键字]
\begin{itemize}
  \item \texttt{try}：包裹可能出错代码。
  \item \texttt{catch}：捕获并处理异常。
  \item \texttt{finally}：通常无论是否异常都会执行。
  \item \texttt{throws}：声明方法可能抛出异常。
  \item \texttt{throw}：主动抛出异常对象。
\end{itemize}
\end{keybox}

\subsection{File 与输入输出流}
\begin{keybox}[区别]
\texttt{File} 表示文件或目录路径本身，不负责读写文件内容。字节流以字节为单位，适合图片、音频等二进制文件；字符流以字符为单位，适合文本文件。读到文件末尾常用返回值 $-1$ 表示。
\end{keybox}

\subsection{线程与 GUI 事件}
\begin{longtable}{p{0.25\textwidth}p{0.65\textwidth}}
\toprule
\textbf{知识点} & \textbf{考试写法} \\
\midrule
启动线程 & 调用 \texttt{start()}，不要直接调用 \texttt{run()} 当作启动新线程 \\
事件源 & 产生事件的组件，如按钮、文本框 \\
事件对象 & 保存事件相关信息 \\
监听器 & 监听事件并执行处理代码 \\
注解 & 给程序元素附加元数据信息，属于了解性考点 \\
\bottomrule
\end{longtable}

\section{程序阅读题模板}

\begin{methodbox}[构造方法阅读]
\begin{enumerate}
  \item 创建对象先执行构造方法。
  \item \texttt{this(...)} 或 \texttt{super(...)} 会先跳转调用另一个构造。
  \item 成员变量没有显式赋值时用默认值。
  \item 最后按 \texttt{System.out.println} 顺序写输出。
\end{enumerate}
\end{methodbox}

\begin{methodbox}[集合输出阅读]
\begin{enumerate}
  \item \texttt{ArrayList} 保留插入顺序，允许重复。
  \item \texttt{HashSet} 不保证顺序，去重。
  \item \texttt{TreeSet} 去重并排序。
  \item 增强 for 按集合迭代顺序输出。
\end{enumerate}
\end{methodbox}

\section{程序设计题模板}

\begin{methodbox}[管理类题固定结构]
\begin{Code}
class Student {
  private String id;
  private String name;
  private double score;
  public Student(String id, String name, double score) {
    this.id = id;
    this.name = name;
    this.score = score;
  }
  public String getId() { return id; }
  public String getName() { return name; }
  public double getScore() { return score; }
}

class ScoreManager {
  private ArrayList<Student> list = new ArrayList<Student>();
  public void addStudent(Student s) { list.add(s); }
  public double avgScore() {
    if (list.isEmpty()) return 0;
    double sum = 0;
    for (Student s : list) sum += s.getScore();
    return sum / list.size();
  }
}
\end{Code}
\end{methodbox}

\begin{warnbox}[程序设计扣分点]
\begin{enumerate}
  \item 忘记导入 \texttt{java.util.*} 或 \texttt{ArrayList}。
  \item 属性没有私有化，getter/setter 不完整。
  \item 方法返回值与题目要求不一致，例如应返回 \texttt{boolean} 却只打印。
  \item 字符串比较用 \texttt{==}，应使用 \texttt{equals()}。
  \item main 方法没有创建对象和测试，导致程序结构不完整。
\end{enumerate}
\end{warnbox}

\section{高频题型方法模板}

\begin{methodbox}[题型一：概念选择填空]
\textbf{识别特征：}题干问关键字、类关系、集合特点、异常语法。\par
\textbf{固定步骤：}先判断属于语法、OOP、集合还是异常；再排除语义明显错误选项。\par
\textbf{易错点：}\texttt{extends} 与 \texttt{implements}、\texttt{throw} 与 \texttt{throws}、\texttt{start()} 与 \texttt{run()}。
\end{methodbox}

\begin{methodbox}[题型二：程序阅读]
\textbf{识别特征：}给一段完整 Java 代码，要求写输出和原因。\par
\textbf{固定步骤：}标出对象实际类型、变量值变化、方法绑定、集合内容变化；逐条写输出。\par
\textbf{易错点：}多态调用看实际对象类型；\texttt{TreeSet} 会排序并去重。
\end{methodbox}

\begin{methodbox}[题型三：程序设计]
\textbf{识别特征：}要求定义实体类和管理类，用 \texttt{ArrayList} 保存对象。\par
\textbf{固定步骤：}实体类私有属性 + 构造 + getter/setter；管理类保存集合 + 增删查改；main 创建对象测试。\par
\textbf{易错点：}题目要求返回值时不要只打印；查找编号必须用 \texttt{equals()}。
\end{methodbox}

\section{速查表}
\begin{longtable}{p{0.25\textwidth}p{0.35\textwidth}p{0.30\textwidth}}
\toprule
\textbf{模块} & \textbf{关键词/写法} & \textbf{使用场景} \\
\midrule
编译运行 & \texttt{javac}、\texttt{java}、JVM & 平台无关性 \\
对象 & \texttt{new}、构造方法、\texttt{this} & 创建和初始化对象 \\
封装 & \texttt{private}+getter/setter & 程序设计题 \\
继承 & \texttt{extends}、\texttt{super} & 子类复用父类 \\
接口 & \texttt{interface}、\texttt{implements} & 定义规范 \\
多态 & 父类引用指向子类对象 & 程序阅读题 \\
类型判断 & \texttt{instanceof} & 向下转型前检查 \\
String & \texttt{equals()}、\texttt{charAt()} & 字符串题 \\
集合 & \texttt{ArrayList}、\texttt{HashSet}、\texttt{TreeSet} & 保存对象、去重、排序 \\
异常 & \texttt{try/catch/finally/throws} & 异常处理 \\
线程 & \texttt{start()} & 启动线程 \\
GUI & 事件源、事件对象、监听器 & 简答题 \\
\bottomrule
\end{longtable}

\section{考前 3～7 天复习计划}
\begin{methodbox}[7 天版]
\begin{enumerate}
  \item Day 1：JDK/JRE/JVM、基本语法、类型转换、数组。
  \item Day 2：类、对象、封装、构造方法、this/static/final。
  \item Day 3：继承、重写、多态、接口、抽象类。
  \item Day 4：String、包装类、集合、泛型。
  \item Day 5：异常、文件流、线程、GUI 事件。
  \item Day 6：集中练程序阅读题，逐行写输出原因。
  \item Day 7：默写管理类程序设计模板，做一套模拟卷并订正。
\end{enumerate}
\end{methodbox}

\begin{methodbox}[3 天版]
\begin{enumerate}
  \item Day 1：背关键字和 OOP 概念，刷选择填空。
  \item Day 2：练继承多态、String、集合输出题。
  \item Day 3：默写实体类 + Manager 类 + main 测试模板。
\end{enumerate}
\end{methodbox}

\section{考前最后 30 分钟速记清单}
\begin{keybox}[只看这些]
\begin{enumerate}
  \item \texttt{.java} 编译成 \texttt{.class}；字节码运行在 JVM 上。
  \item public 类名必须与源文件名一致；入口方法名是 \texttt{main}。
  \item 构造方法无返回值、可重载、不能继承；\texttt{this(...)} 和 \texttt{super(...)} 必须第一句。
  \item 封装：\texttt{private} 属性 + public 方法访问。
  \item 继承用 \texttt{extends}；实现接口用 \texttt{implements}。
  \item 多态：父类引用指向子类对象，重写方法执行子类版本。
  \item 向下转型前用 \texttt{instanceof}；字符串内容比较用 \texttt{equals()}。
  \item \texttt{ArrayList} 有序可重复；\texttt{TreeSet} 排序去重。
  \item 异常：\texttt{catch} 捕获，\texttt{finally} 通常必执行，\texttt{throws} 声明可能抛出。
  \item 线程启动用 \texttt{start()}；程序设计题先保证类结构完整。
\end{enumerate}
\end{keybox}

\end{document}
